% Part: incompleteness
% Chapter: arithmetization-syntax
% Section: coding-terms

\documentclass[../../../include/open-logic-section]{subfiles}

\begin{document}

\olfileid[hi]{inc}{art}{trm}
\olsection{पदों का कूटलेखन}

\begin{explain}
पद केवल एक विशेष प्रकार का चिह्न-अनुक्रम है: पदों के निर्माण-नियमों के अनुसार
वह अचरों और चरों से आगमनात्मक रूप में बनता है। चूँकि चिह्नों की कूटलेखन-योजना
और संख्या-अनुक्रमों को कूटित करने के किसी उपाय से चिह्न-अनुक्रमों को संख्याओं
के रूप में कूटित किया जा सकता है, इसलिए पदों को ग्योडेल (G\"odel) संख्याएँ
देना कठिन नहीं। चुनौती यह दिखाना है कि किसी संख्या का किसी सुगठित पद की
ग्योडेल (G\"odel) संख्या होना संगणनीय, बल्कि आदिम पुनरावर्ती, गुण है।
\end{explain}

!!^{variable} और !!{constant} सबसे सरल पद हैं, और यह जाँचना सरल है कि $x$
ऐसे किसी पद की ग्योडेल (G\"odel) संख्या है या नहीं: $\fn{Var}(x)$ तब लागू
होता है जब किसी~$i$ के लिए $x$, $\Gn{\Obj v_i}$ हो। दूसरे शब्दों में, $x$
लंबाई~$1$ वाला अनुक्रम है और उसका एकमात्र अवयव $(x)_0$ किसी
!!{variable}~$\Obj v_i$ का कूट है; अर्थात् किसी~$i$ के लिए $x$,
$\tuple{\tuple{1, i}}$ है। इसी प्रकार, $\fn{Const}(x)$ तब लागू होता है जब
किसी~$i$ के लिए $x$, $\Gn{\Obj c_i}$ हो। ये दोनों संबंध आदिम पुनरावर्ती हैं,
क्योंकि यदि ऐसा $i$ विद्यमान है, तो वह अनिवार्यतः $< x$ है:
\begin{align*}
  \fn{Var}(x) & \defiff \bexists{i<x}{x = \tuple{\tuple{1, i}}}\\
  \fn{Const}(x) & \defiff \bexists{i<x}{x = \tuple{\tuple{2, i}}}
\end{align*}

\begin{prop}
\ollabel{prop:term-primrec}
संबंध $\fn{Term}(x)$ और $\fn{ClTerm}(x)$, जो क्रमशः तभी लागू होते हैं जब $x$
किसी पद या बंद पद की ग्योडेल (G\"odel) संख्या हो, आदिम पुनरावर्ती हैं।
\end{prop}

\begin{proof}
चिह्नों का अनुक्रम~$s$ पद है यदि और केवल यदि पदों का कोई अनुक्रम~$s_0$,
\dots, $s_{k-1} = s$ यह दर्ज करता है कि पदों के निर्माण-नियमों के अनुसार
!!{constant} और !!{variable} से पद~$s$ कैसे बना। यह व्यक्त करने के लिए कि ऐसा
अभिकल्पित निर्माण-अनुक्रम निर्माण-नियमों का पालन करता है, प्रत्येक $i < k$ के
लिए निम्न में से कोई एक बात सत्य होनी चाहिए:
\begin{enumerate}
\item $s_i$, !!a{variable}~$\Obj v_j$ है; अथवा
\item $s_i$, !!a{constant}~$\Obj c_j$ है; अथवा
\item $s_i$, स्थान~$i$ से पहले आने वाले $n$ पदों $t_1$, \dots, $t_n$ से
  किसी $n$-स्थानी !!{function}~$\Obj f^n_j$ का उपयोग करके बना है।
\end{enumerate}
यह दिखाने के लिए कि ग्योडेल (G\"odel) संख्याओं पर संगत संबंध आदिम पुनरावर्ती
है, हमें इस शर्त को आदिम पुनरावर्ती ढंग से, अर्थात् आदिम पुनरावर्ती फलनों,
संबंधों और परिबद्ध परिमाणीकरण का उपयोग करके, व्यक्त करना होगा।

मान लें कि $y$ वह संख्या है जो अनुक्रम $s_0$, \dots, $s_{k-1}$ को कूटित करती
है; अर्थात् $y = \tuple{\Gn{s_0}, \dots, \Gn{s_{k-1}}}$। वह ग्योडेल (G\"odel)
 संख्या~$x$ वाले पद के निर्माण-अनुक्रम को कूटित करती है यदि और केवल
यदि सभी $i < k$ के लिए:
\begin{enumerate}
\item $\fn{Var}((y)_i)$, अथवा
\item $\fn{Const}((y)_i)$, अथवा
\item कोई $n$ और संख्या~$z = \tuple{z_1, \dots, z_n}$ ऐसी है कि प्रत्येक
  $z_l$, किसी $i' < i$ के लिए $(y)_{i'}$ के बराबर है, और
\[
(y)_i = \Gn{\Obj f^n_j(} \concat \fn{flatten}(z) \concat \Gn{)},
\]
\end{enumerate}
और इसके अतिरिक्त $(y)_{k-1} = x$। (फलन $\fn{flatten}(z)$ अनुक्रम
$\tuple{\Gn{t_1}, \dots, \Gn{t_n}}$ को $\Gn{t_1, \dots, t_n}$ में बदलता है
और आदिम पुनरावर्ती है।)

खंड (3) में सूचकांक $j$, $n$, पदों $t_l$ की ग्योडेल (G\"odel) संख्याएँ
$z_l$, और अनुक्रम~$\tuple{z_1, \dots, z_n}$ का कूट~$z$, सभी~$y$ से छोटे
हैं। ऊपर $k$ के स्थान पर $\len{y}$ रख सकते हैं। अतः हम ``$y$, ग्योडेल (G\"odel)
 संख्या~$x$ वाले पद के निर्माण-अनुक्रम का कूट है'' को ऐसे ढंग से
व्यक्त कर सकते हैं जिससे यह संबंध आदिम पुनरावर्ती सिद्ध होता है।

अब हमें केवल यह विश्वास दिलाना है कि~$y$ पर कोई आदिम पुनरावर्ती परिसीमा है।
किंतु यदि $x$ किसी पद की ग्योडेल (G\"odel) संख्या है, तो उसका ऐसा
निर्माण-अनुक्रम होना चाहिए जिसमें अधिकतम $\len{x}$ पद हों (क्योंकि~$s$ के
निर्माण-अनुक्रम का प्रत्येक पद~$s$ में किसी स्थान से आरंभ होना चाहिए, और दो
उपपद एक ही स्थान से आरंभ नहीं हो सकते)। निस्संदेह,~$s$ के प्रत्येक उपपद की
ग्योडेल (G\"odel) संख्या $\le x$ है। अतः सदा कोई ऐसा निर्माण-अनुक्रम है
जिसका कूट $\le p_{k-1}^{k(x+1)}$ है, जहाँ $k=\len{x}$।

$\fn{ClTerm}$ के लिए केवल !!{variable} वाला खंड छोड़ दें।
\end{proof}

\begin{prob}
दिखाएँ कि अनुक्रम $\tuple{\Gn{t_1}, \dots, \Gn{t_n}}$ को $\Gn{t_1, \dots,
t_n}$ में बदलने वाला फलन $\fn{flatten}(z)$ आदिम पुनरावर्ती है।
\end{prob}

\begin{prop}\ollabel{prop:num-primrec}
  फलन $\fn{num}(n) = \Gn{\num{n}}$ आदिम पुनरावर्ती है।
\end{prop}

\begin{proof}
  हम $\fn{num}(n)$ को आदिम पुनरावर्तन द्वारा परिभाषित करते हैं:
  \begin{align*}
    \fn{num}(0) & = \Gn{\Obj 0}\\
    \fn{num}(n+1) & = \Gn{\prime(} \concat \fn{num}(n) \concat \Gn{)}.
  \end{align*}
\end{proof}

\end{document}
